\section{Reproducibility, availability and limitations}\label{sec:repro}

\subsection{Identified private software snapshot}

The accompanying archive \path{pyexpint_review_package_20261011.zip}
contains the software, original experiment records, distribution artifacts
and validation logs.  Its source tree, excluding the packaging-only
\texttt{SNAPSHOT.txt}, has the Git tree identifier
\begin{center}\small\ttfamily
9f62cbf6a6455868e358386ccff824507287aeb3.
\end{center}
This identifier can be reconstructed from the supplied contents and file
modes.  The recorded preparation commit is
\texttt{0ad4b3c4786086a04cf22be33375cc9997ff1344}.  The manifest covers
317 payload files.  Both distributions contain package version metadata
\releaseversion{}, retained from development; this label is not the identity
of the corrected snapshot and is not a claim that it is the earlier public
release.  The 20 Python source files in the supplied wheel match those in
\path{software/src/pyexpint}.  The source distribution likewise matches the
corresponding supplied source files apart from generated build metadata.
No new public release is implied here.  The earlier public 0.9.0 archive,
DOI \path{10.5281/zenodo.22820595}~\cite{VillatoroPyEXPINT090}, remains a
separate historical artifact.  Software licensing is MIT and original
non-software content is CC BY~4.0; original EXPINT MATLAB sources are not
redistributed.

The detailed preparation-environment JUnit records report 275 passing tests
from the snapshot and the same 275 identifiers against the installed wheel.
The logged environment is Python 3.11.15, NumPy 1.24.3, SciPy 1.10.1 and
pytest 7.1.2.  Four examples ran outside the repository; the source
distribution was also built and installed and ran the no-exact-history
example.  Scripts print a fixed ``WSL'' label, but the package explicitly
identifies these detailed logs as preparation-container runs.  An additional
author-workstation summary concerns Python 3.11.16 and earlier commits; it
is not substituted for those detailed logs.  A Python 3.11--3.14 workflow is
configured in the source, but execution of that entire matrix is not
established by the supplied candidate records.

\subsection{Experiment and manuscript reconstruction}

Table~\ref{tab:protocol} distinguishes original records, candidate
comparisons and historical results.  The full original Allen--Cahn JSON
replaces the provisional transcription and records 90 individual times,
step sequences, construction coefficients and initial/reference states.
Its source identifier and environment remain those of the workstation run.
The driver and original records are in the private companion, and the JSON
used by Figure~\ref{fig:broadband} is byte-identical to that original.
Reference-comparison scalars and per-solve error scalars retain their status
as driver outputs when both compared final states are not archived.

The 72 original canonical JSON files are now supplied under
\path{software/reproducibility/toms/rep1/raw} and the corresponding
\texttt{rep2} and \texttt{rep3} directories.  They contain 864 complete-solve
timing records.  Matching block identifiers produces the 54 paired
summaries used in this article.  Those times remain measurements of the
historical implementation, identified by its retained source manifest.
The candidate's \path{tools/differential_canonical.py} has a separate log
with 72 comparisons against the three campaigns and no recorded
mismatches under its stated error/count criteria.  Its accepted-step
sequences are outputs of the candidate run, not reconstructed historical
sequences.  This check neither establishes equal runtime nor reexecutes the
108-case fixed-step method campaign, which remains historical evidence.
No new cluster timing of the private snapshot is claimed.

The cross-check companion provides the two drivers, 102 original/corrected
scheme output rows, coefficient values, refinement output, correction diffs
and selected original-file hashes.  Supplementary Table~S6 still reproduces
the unchanged summary; the additional files make its test configuration and
source boundaries inspectable.  The supplied coverage index connects all
47 catalogue entries with source definitions and concrete test identifiers;
it distinguishes structural, coefficient, scalar-order and linear tests
from declared order metadata.

For manuscript reconstruction,
\path{figure_sources/make_figures.py} regenerates the numerical panels and
shared legend from the selected records.
\path{tables/make_tables.py} generates four tables (startup, catalogue,
Lawson and development step doubling), and
\path{tables/make_crosscheck_table.py} generates Table~S6 from the preserved
summary.  Other tables are supplied LaTeX compositions.  The figure and
table transformations are separate from integrator execution and from the
recorded installation tests.  An evidence index identifies the private
archive and the paths of its logs without replacing historical identifiers.

\subsection{AI-assisted research and manuscript workflow}\label{sec:ai-use}

Generative-AI tools were used during the research lifecycle as interactive
assistants: ChatGPT (OpenAI) for proposing diagnostic experiments, drafting and
checking shell and Python snippets, organizing software-release audits and
assembling an earlier draft of this manuscript, and Claude (Anthropic, model
Claude Fable 5.1) for the independent review of that draft, the correction
blocks of the software, the cross-check against the original EXPINT sources
and the revision of the text.  The tools did not generate the numerical values
reported as experimental results; those come from archived \pkg{} executions
and campaign records.  Checks described in the manuscript include code execution, regression
tests, comparison of preserved sources and algebraic deductions.  These
forms of evidence are identified separately; the algebraic ETD34 example
is not represented as a new numerical execution, and execution logs are distinguished from AI-generated descriptions.  The experimental scope, the scientific
interpretation, the code admitted to the release and the claims of the
manuscript are the responsibility of the human author.  This section is
included because the assistance covered experiment design, coding, analysis,
validation and reproducibility, not only editing.

\subsection{Threats to validity and claim boundaries}

\paragraph{Wall time is local evidence.}
The canonical timing campaigns are replicated executions on a specified
cluster environment, not a cross-machine performance study.  Two of the three
ran on the same node, and the between-replica spread of the paired ratios is
several times their within-replica interquartile range.  The paper uses timing
ratios as a demonstration of the benchmarking infrastructure and relies on
achieved error and work counters for portable interpretation.

\paragraph{The canonical problem is a single Fourier mode.}
Its manufactured solution is an eigenvector of the periodic operator, so the
vectors acted upon are dominated by a few modes and the Krylov counts depend
weakly on the grid in the measured cases.  The campaign says nothing about
operators with broadband data, which the problem of Section~\ref{sec:broadband}
addresses through the supplied original workstation record; half of the
canonical campaign is maximum-step-limited.

\paragraph{A work counter is not total cost.}
A matrix-vector product in Arnoldi, Leja or a hybrid path has different
associated vector and projected-matrix work, which is why the Leja backend
costs 1.08 to 1.47 times the KIOPS-type time with 1.63 to 3.11 times its
products.  The hybrid counter omits the internal SciPy products.  Estimated
algorithmic workspace is not peak resident memory.

\paragraph{The public Leja path has no substepping.}
The real-Leja backend interpolates on the whole scaled spectral interval with
a maximum degree of 80: degree 69 is used at width 240 on a dissipative
operator and the cap is reached near width 300 at tolerance $10^{-10}$; the
canonical problem reaches width 53 and the broadband problem 210.  Wider
intervals require a larger degree or a smaller step, and the failure is
reported as a rejected step by the adaptive solvers.

\paragraph{Tolerance semantics are not automatically equivalent.}
Temporal tolerances, matrix-action tolerances and stopping criteria play
different roles across realizations.  The canonical timing campaign records an
action fraction of 0.01; the package default is 0.05.  The work--precision
experiments either fix the action backend and tolerance or display achieved
error, rather than treating a requested tolerance as an accuracy identity.

\paragraph{The ETD34 estimator is blind to state-independent forcing.}
Section~\ref{sec:adaptivity} documents that the published pair estimates no
error for $N(t,y)=g(t)$, and the adaptive comparisons of
Section~\ref{sec:adaptivework} are read with that confound in mind.

\paragraph{Exact-history experiments isolate method bodies.}
The fixed-window method work--precision experiment uses exact nonlinear
history before the common numerical interval.  This isolates the body of
multistep and general linear schemes; with the default numerical startup the
guaranteed end-to-end order is $\min(p,5)$ and the observed slopes for the selected
order six and seven entries between $h=0.05$ and $h=0.025$ lie between
4.55 and 4.80
(Section~\ref{sec:startup}).

\paragraph{Automatic selection is heuristic.}
The rule-based selector uses operator structure, problem size and a scaled
spectral-width threshold.  On the canonical problem it chose the slower
backend in every configuration.  Its thresholds are exposed, a calibrated
selector interface exists, and the current evidence does not support a
globally optimal selector.

\paragraph{No external-speed or distributed-scaling claim.}
The relationship to other packages is established functionally in
Section~\ref{sec:related} and, for the historical catalogue, by execution of
the original EXPINT sources; no controlled runtime comparison against external
packages is part of this article, and the release has no distributed backend
contract.
